\documentclass{article}
\usepackage{/globals/tables}
\usepackage{/globals/de}
\usepackage{/globals/math}
\usepackage{/globals/anki}
\usepackage{hyperref}
\ankiprefix{aussagenlogik}
\ankitopic{Aussagenlogik}
\ankishow
\begin{document}
\section{Aussagenlogik} 
Die Aussagenlogik beschäftigt sich mit Aussagen. Eine Aussage ist ein Satz der wahr oder falsch sein kann, hier dargestellt durch jeweils eine 0 und eine 1.

Aussagen können zusammensetzungen aus mehreren anderen Aussagen sein, oder sie können \emph{atomar} sein.
\ankinote{arten}{Aussagen Arten}{atomar, zusammengesetzt}

\subsection{Verknüpfungen}
\begin{description}
 \item[$\neg A$] Die \emph{Negation} \textquote{nicht A} ist wahr wenn A nicht wahr ist
 \item[$A \land B$] Die \emph{Konjunktion} \textquote{A und B} ist wahr wenn sowohl A wahr ist und B wahr ist.
 \item[$A \lor B$] Die \emph{Disjunktion} \textquote{A oder B} ist wahr wenn mindestens eine der beiden Formeln wahr ist.
 \item[$A \oplus B$] Die \emph{Kontravalenz} \textquote{entweder A oder B} ist wahr wenn genau eine der beiden Formeln wahr ist. 
 \item[$A \implies B$] Die \emph{Implikation} \textquote{wenn A dann B} ist wahr wenn A falsch ist oder wenn A wahr ist und B wahr ist.
 \item[$A \iff B$] Die \emph{Äquivalenz} \textquote{A genau dann wenn B} ist wahr wenn A und B den gleichen Wert haben.
\end{description}
\ankinote{not}{\(\neg A\)}{Negation, nicht A}
\ankinote{and}{\(A \land B\)}{Konjunktion, A und B}
\ankinote{or}{\(A \lor B\)}{Disjunktion, A oder B}
\ankinote{xor}{\(A \oplus B\)}{Kontravalenz, A xor B}
\ankinote{implies}{\(A \implies B\)}{Implikation, wenn A dann B}
\ankinote{iff}{\(A \iff B\)}{Äquivalenz, A genau dann wenn B}

\subsection{Sprache Formalisieren}
Unsere menschliche Sprache ist nicht formal und oftmals auch ohne weiterem Vorwissen sehr ungenau und missverständlich.

Sie kann aber formalisiert werden indem alle atomaren Aussagensymbole einer Aussage gefunden werden, gemäß des Satzes zu Teilaussagen kombiniert werden, die dann zu einer Gesamtaussage kombiniert werden können.

\subsection{Normalformen}
Eine Aussage ist ein der Normalform wenn sie nur aus Konjunktionen, Disjunktionen und Negationen besteht. Genauer gibt es zwei Arten an Normalformen.
\begin{description}
 \item[DNF] Eine disjunktive Normalform (DNF) besteht nur aus Disjunktionen von Konjunktionen.
 \item[KNF] Eine konjunktive Normalform (DNF) besteht nur aus Konjunktionen von Disjunktionen.
\end{description}
Wie DNFs und KNFs gebildet werden können ist weiter in \hyperref[sec:normalform_bilden]{Einführung Informatik} beschrieben.

\subsection{Wahrheitstabellen}
Eine Wahrheitstabelle ist eine tabellarische Darstellung aller möglichen Kombinationen der atomaren Formeln einer Formeln und den dazugehörigen Werten der Formel.
\begin{ztable}{cc|c}{$A$ & $B$ & $A \land B$}
 $0$ & $0$ & $0$ \\
 $0$ & $1$ & $0$ \\
 $1$ & $0$ & $0$ \\
 $1$ & $1$ & $1$ \\
\end{ztable}

\subsection{Tautologien und Kontradiktionen}
\emph{Tautologien} und \emph{Kontradiktionen} sind besondere Arten an Formeln; ihre Ausgabe hängt nicht von den Eingaben ab.

Eine Tautologie ist eine Aussage die \emph{immer}, unabhängig von der Wahrheit der Eingaben, selbst wahr ist.
\ankinote{taut}{Tautologie}{immer wahr}

Eine Kontradiktion ist gegenteilig immer falsch.\ankinote{contradict}{Kontradiktion}{immer falsch}
\end{document}